% v2.6.0 figure UID: FIG-P067-01
% Image-reference reverse redraw: exact CDF/PMF coordinates retained; paired reading and marker clarity refined.
\tikzset{
  slfig-FIG-P067-01/.style={every path/.append style={line cap=round,line join=round}},
  slfig-FIG-P067-01-guide/.style={
    draw=SLGray!78,line width=.50pt,dash pattern=on 2.4pt off 1.7pt
  },
  slfig-FIG-P067-01-mass/.style={
    anchor=west,font=\fontsize{9.2pt}{11.0pt}\selectfont,
    fill=white,fill opacity=.95,text opacity=1,inner sep=.9pt
  },
  slfig-FIG-P067-01-note/.style={
    font=\fontsize{9.2pt}{11.2pt}\selectfont,text=SLInk!82,
    fill=white,fill opacity=.94,text opacity=1,inner sep=.8pt
  }
}
\pgfplotsset{slfig-FIG-P067-01-axis/.style={
  axis line style={draw=SLInk!90,line width=.60pt},
  tick style={draw=SLInk!80,line width=.55pt},
  tick label style={font=\fontsize{8.8pt}{10.5pt}\selectfont},
  label style={font=\fontsize{9.4pt}{11.2pt}\selectfont},
  clip=false
}}
\pgfkeysifdefined{/tikz/alt}{}{\tikzset{alt/.code={}}}
\begin{figure}[H]
\centering
\begin{tikzpicture}[slfig-FIG-P067-01,alt={对象—关系—结论：离散随机变量的分布函数：跳跃高度等于对应点的概率质量}]
\begin{axis}[slfig-FIG-P067-01-axis,
  name=cdf,width=.82\linewidth,height=3.5cm,
  axis lines=left,grid=none,clip=false,
  xmin=.45,xmax=4.55,ymin=-.03,ymax=1.08,
  xtick={1,2,3,4},xticklabels={,,,},
  ytick={0,.45,.8,1},ylabel={$F_X(t)$},
  tick label style={font=\fontsize{8.8pt}{10.5pt}\selectfont},
  label style={font=\fontsize{9.4pt}{11.2pt}\selectfont}
]
  \addplot[const plot mark left,draw=SLBlue,line width=1.00pt,no marks]
    coordinates {(.5,0) (1,.15) (2,.45) (3,.80) (4,1) (4.5,1)};
  \addplot[only marks,mark=*,mark size=2.2pt,
    mark options={draw=SLBlue,fill=SLBlue,line width=.70pt}]
    coordinates {(1,.15) (2,.45) (3,.80) (4,1)};
  \addplot[only marks,mark=o,mark size=2.35pt,
    mark options={draw=SLBlue,fill=white,line width=.75pt}]
    coordinates {(1,0) (2,.15) (3,.45) (4,.80)};
  \addplot[draw=SLGray!82,line width=.55pt,dash pattern=on 3pt off 2pt,no marks]
    coordinates {(.5,1) (4.5,1)};
  \draw[slfig-FIG-P067-01-guide] (axis cs:1,0)--(axis cs:1,1.03);
  \draw[slfig-FIG-P067-01-guide] (axis cs:2,0)--(axis cs:2,1.03);
  \draw[slfig-FIG-P067-01-guide] (axis cs:3,0)--(axis cs:3,1.03);
  \draw[slfig-FIG-P067-01-guide] (axis cs:4,0)--(axis cs:4,1.03);
  \node[slfig-FIG-P067-01-mass]
    at (axis cs:1.08,.09) {$p_1$};
  \node[slfig-FIG-P067-01-mass]
    at (axis cs:2.06,.30) {$p_2$};
  \node[slfig-FIG-P067-01-mass]
    at (axis cs:3.06,.625) {$p_3$};
  \node[slfig-FIG-P067-01-mass]
    at (axis cs:4.08,.85) {$p_4$};
  \node[slfig-FIG-P067-01-note,fill=none,anchor=north west]
    at (axis cs:.58,.87) {右连续：实心点取跳后值};

\end{axis}
\begin{axis}[
  at={(cdf.south west)},anchor=north west,yshift=-3mm,
  width=.82\linewidth,height=2.9cm,
  axis lines=left,grid=none,clip=false,
  xmin=.45,xmax=4.55,ymin=0,ymax=.41,
  xtick={1,2,3,4},ytick={0,.15,.30,.35},yticklabels={0,0.15,,0.35},
  xlabel={$t$},ylabel={$p_X(t)$},
  tick label style={font=\fontsize{8.6pt}{10.3pt}\selectfont},
  label style={font=\fontsize{9.4pt}{11.2pt}\selectfont}
]
  \addplot[ycomb,draw=SLTeal,line width=1.00pt,mark=*,mark size=2.2pt,
    mark options={draw=SLTeal,fill=SLInk,line width=.65pt}]
    coordinates {(1,.15) (2,.30) (3,.35) (4,.20)};
  \node[anchor=east,xshift=-2pt,yshift=-4.5pt,font=\fontsize{8.6pt}{10.3pt}\selectfont]
    at (axis cs:.45,.30) {$0.3$};
  \draw[slfig-FIG-P067-01-guide] (axis cs:1,0)--(axis cs:1,.40);
  \draw[slfig-FIG-P067-01-guide] (axis cs:2,0)--(axis cs:2,.40);
  \draw[slfig-FIG-P067-01-guide] (axis cs:3,0)--(axis cs:3,.40);
  \draw[slfig-FIG-P067-01-guide] (axis cs:4,0)--(axis cs:4,.40);
  \node[slfig-FIG-P067-01-note,anchor=north east]
    at (axis cs:4.45,.39) {同一 $t_i$：跳高 $=p_i$};
\end{axis}
\end{tikzpicture}
\caption{离散随机变量的分布函数：跳跃高度等于对应点的概率质量}\label{fig:V1-C04-cdf}
\end{figure}
